% GENERATED FILE. Edit canonical lessons, then run npm run generate:books.
% canonical-source-hash: fnv1a64:6ba0198de961a614
% canonical-lessons: KA-C403-samaja, KA-C403-svavalambane, KA-C403-jagrti, KA-C403-vaividhya, KA-C403-rupisu

\chapter{Society and Self-Reliance}
\label{ch:ka-society-and-self-reliance}
\hlchaptermodality{403}

\begin{chapteropening}
\textbf{By the end of this chapter:} \emph{I can say society, self-reliance, awareness, diversity and to devise.}

Complete the last lesson of chapter 403: I can say society, self-reliance, awareness, diversity and to devise.
\end{chapteropening}

\section[\texorpdfstring{\kn{ಸಮಾಜ} (samāja)}{ಸಮಾಜ (samāja)}]{\kn{ಸಮಾಜ} (samāja) — society}

\label{lesson:KA-C403-samaja}



\begin{quote}
Before the new one: say the Kannada for to conceal, then the Kannada for to recommend.
\end{quote}

\subsection*{You'll want to know: \kn{ಸಮಾಜ}}
\textbf{\kn{ಸಮಾಜ}} — \emph{samāja} — \textquotedblleft{}society\textquotedblright{}.

\begin{quote}
\kn{ಸಮಾಜದಲ್ಲಿ} \kn{ಎಲ್ಲರಿಗೂ} \kn{ಗೌರವ} \kn{ಸಿಗಬೇಕು}.
\end{quote}

\begin{quote}
\emph{samājadalli ellarigū gaurava sigabēku.}
\end{quote}

\begin{quote}
Everyone in society should be respected.
\end{quote}

\kn{ಸಾಮಾಜಿಕ}, social.

\begin{cousinweb}
\textbf{\kn{ಸಮಾಜ}} (\emph{samāja}) is borrowed from Sanskrit \emph{samāja}, \textquotedblleft{}a meeting, an assembly\textquotedblright{} (CDIAL 13207).
\end{cousinweb}

\begin{grammarlens}[title={What you've built}]
One more word for giving reasons and opinions.
\end{grammarlens}

\subsection*{Your turn}
\begin{itemize}
\raggedright
  \item \emph{Say it:} \emph{samāja}
  \item \emph{Say it:} \emph{samāja}, once more
  \item \emph{Say it:} \emph{samājadalli ellarigū gaurava sigabēku.} — Everyone in society should be respected.
  \item \emph{Recall:} say the Kannada for to conceal, then the Kannada for to recommend, then say \emph{samāja} again
\end{itemize}

\begin{tcolorbox}[breakable,colback=teal!4,colframe=teal!35!black,title={Before you move on}]
What does \kn{ಸಮಾಜ} mean? (\textquotedblleft{}Society\textquotedblright{}.) Say it once more.
\end{tcolorbox}
\section[\texorpdfstring{\kn{ಸ್ವಾವಲಂಬನೆ} (svāvalambane)}{ಸ್ವಾವಲಂಬನೆ (svāvalambane)}]{\kn{ಸ್ವಾವಲಂಬನೆ} (svāvalambane) — self-reliance}

\label{lesson:KA-C403-svavalambane}



\begin{quote}
Before the new one: say the Kannada for to recommend, then the Kannada for society.
\end{quote}

\subsection*{You'll want to know: \kn{ಸ್ವಾವಲಂಬನೆ}}
\textbf{\kn{ಸ್ವಾವಲಂಬನೆ}} — \emph{svāvalambane} — \textquotedblleft{}self-reliance\textquotedblright{}.

\begin{quote}
\kn{ಹಳ್ಳಿಯ} \kn{ಹೆಂಗಸರು} \kn{ಹೊಲಿಗೆ} \kn{ಕಲಿತು} \kn{ಸ್ವಾವಲಂಬನೆ} \kn{ಸಾಧಿಸಿದರು}.
\end{quote}

\begin{quote}
\emph{haḷḷiya heṅgasaru holige kalitu svāvalambane sādhisidaru.}
\end{quote}

\begin{quote}
The village women learned tailoring and became self-reliant.
\end{quote}

\kn{ಸ್ವ} (self) + \kn{ಅವಲಂಬನೆ} (leaning on).

\begin{cousinweb}
\textbf{\kn{ಸ್ವಾವಲಂಬನೆ}} (\emph{svāvalambane}) is built from Sanskrit \emph{sva}, \textquotedblleft{}self\textquotedblright{}, and \emph{avalambana}, \textquotedblleft{}leaning on, depending upon\textquotedblright{} (Monier-Williams): depending on oneself.
\end{cousinweb}

\begin{grammarlens}[title={What you've built}]
One more word for giving reasons and opinions.
\end{grammarlens}

\subsection*{Your turn}
\begin{itemize}
\raggedright
  \item \emph{Say it:} \emph{svāvalambane}
  \item \emph{Say it:} \emph{svāvalambane}, once more
  \item \emph{Say it:} \emph{haḷḷiya heṅgasaru holige kalitu svāvalambane sādhisidaru.} — The village women learned tailoring and became self-reliant.
  \item \emph{Recall:} say the Kannada for to recommend, then the Kannada for society, then say \emph{svāvalambane} again
\end{itemize}

\begin{tcolorbox}[breakable,colback=teal!4,colframe=teal!35!black,title={Before you move on}]
What does \kn{ಸ್ವಾವಲಂಬನೆ} mean? (\textquotedblleft{}Self-reliance\textquotedblright{}.) Say it once more.
\end{tcolorbox}
\section[\texorpdfstring{\kn{ಜಾಗೃತಿ} (jāgṛti)}{ಜಾಗೃತಿ (jāgṛti)}]{\kn{ಜಾಗೃತಿ} (jāgṛti) — awareness}

\label{lesson:KA-C403-jagrti}



\begin{quote}
Before the new one: say the Kannada for society, then the Kannada for self-reliance.
\end{quote}

\subsection*{You'll want to know: \kn{ಜಾಗೃತಿ}}
\textbf{\kn{ಜಾಗೃತಿ}} — \emph{jāgṛti} — \textquotedblleft{}awareness\textquotedblright{}.

\begin{quote}
\kn{ಆರೋಗ್ಯದ} \kn{ಬಗ್ಗೆ} \kn{ಜನರಲ್ಲಿ} \kn{ಜಾಗೃತಿ} \kn{ಬೇಕು}.
\end{quote}

\begin{quote}
\emph{ārōgyada bagge janaralli jāgṛti bēku.}
\end{quote}

\begin{quote}
People need awareness about health.
\end{quote}

\kn{ಜಾಗೃತಿ} \kn{ಮೂಡಿಸು}, to raise awareness.

\begin{cousinweb}
\textbf{\kn{ಜಾಗೃತಿ}} (\emph{jāgṛti}) is built on Sanskrit \emph{jāgṛ}, \textquotedblleft{}to be awake, to watch\textquotedblright{} (Monier-Williams): a waking up.
\end{cousinweb}

\begin{grammarlens}[title={What you've built}]
One more word for giving reasons and opinions.
\end{grammarlens}

\subsection*{Your turn}
\begin{itemize}
\raggedright
  \item \emph{Say it:} \emph{jāgṛti}
  \item \emph{Say it:} \emph{jāgṛti}, once more
  \item \emph{Say it:} \emph{ārōgyada bagge janaralli jāgṛti bēku.} — People need awareness about health.
  \item \emph{Recall:} say the Kannada for society, then the Kannada for self-reliance, then say \emph{jāgṛti} again
\end{itemize}

\begin{tcolorbox}[breakable,colback=teal!4,colframe=teal!35!black,title={Before you move on}]
What does \kn{ಜಾಗೃತಿ} mean? (\textquotedblleft{}Awareness\textquotedblright{}.) Say it once more.
\end{tcolorbox}
\section[\texorpdfstring{\kn{ವೈವಿಧ್ಯ} (vaividhya)}{ವೈವಿಧ್ಯ (vaividhya)}]{\kn{ವೈವಿಧ್ಯ} (vaividhya) — diversity; variety}

\label{lesson:KA-C403-vaividhya}



\begin{quote}
Before the new one: say the Kannada for self-reliance, then the Kannada for awareness.
\end{quote}

\subsection*{You'll want to know: \kn{ವೈವಿಧ್ಯ}}
\textbf{\kn{ವೈವಿಧ್ಯ}} — \emph{vaividhya} — \textquotedblleft{}diversity; variety\textquotedblright{}.

\begin{quote}
\kn{ಈ} \kn{ಊರಿನಲ್ಲಿ} \kn{ಭಾಷೆಗಳ} \kn{ವೈವಿಧ್ಯ} \kn{ಇದೆ}.
\end{quote}

\begin{quote}
\emph{ī ūrinalli bhāṣegaḷa vaividhya ide.}
\end{quote}

\begin{quote}
This town has a variety of languages.
\end{quote}

\kn{ವೈವಿಧ್ಯಮಯ}, varied.

\begin{cousinweb}
\textbf{\kn{ವೈವಿಧ್ಯ}} (\emph{vaividhya}) is built on Sanskrit \emph{vividha}, \textquotedblleft{}of various sorts, manifold\textquotedblright{} (Monier-Williams), with the vowel strengthening that makes an abstract noun.
\end{cousinweb}

\begin{grammarlens}[title={What you've built}]
One more word for giving reasons and opinions.
\end{grammarlens}

\subsection*{Your turn}
\begin{itemize}
\raggedright
  \item \emph{Say it:} \emph{vaividhya}
  \item \emph{Say it:} \emph{vaividhya}, once more
  \item \emph{Say it:} \emph{ī ūrinalli bhāṣegaḷa vaividhya ide.} — This town has a variety of languages.
  \item \emph{Recall:} say the Kannada for self-reliance, then the Kannada for awareness, then say \emph{vaividhya} again
\end{itemize}

\begin{tcolorbox}[breakable,colback=teal!4,colframe=teal!35!black,title={Before you move on}]
What does \kn{ವೈವಿಧ್ಯ} mean? (\textquotedblleft{}Diversity; variety\textquotedblright{}.) Say it once more.
\end{tcolorbox}
\section[\texorpdfstring{\kn{ರೂಪಿಸು} (rūpisu)}{ರೂಪಿಸು (rūpisu)}]{\kn{ರೂಪಿಸು} (rūpisu) — to devise; to shape}

\label{lesson:KA-C403-rupisu}



\begin{quote}
Before the new one: say the Kannada for awareness, then the Kannada for diversity.
\end{quote}

\subsection*{You'll want to know: \kn{ರೂಪಿಸು}}
\textbf{\kn{ರೂಪಿಸು}} — \emph{rūpisu} — \textquotedblleft{}to devise; to shape\textquotedblright{}.

\begin{quote}
\kn{ಶಿಕ್ಷಕರು} \kn{ಹೊಸ} \kn{ಪಾಠಗಳನ್ನು} \kn{ರೂಪಿಸಿದರು}.
\end{quote}

\begin{quote}
\emph{śikṣakaru hosa pāṭhagaḷannu rūpisidaru.}
\end{quote}

\begin{quote}
The teachers devised new lessons.
\end{quote}

Past (I …): \kn{ರೂಪಿಸಿದೆ}, \emph{rūpiside}.

What is shaped takes -\kn{ಅನ್ನು}.

\begin{cousinweb}
\textbf{\kn{ರೂಪಿಸು}} (\emph{rūpisu}) is built on Sanskrit \emph{rūpa}, \textquotedblleft{}form, shape\textquotedblright{} (CDIAL 10803), with the Kannada verb ending -\emph{isu}: to give form to.
\end{cousinweb}

\begin{grammarlens}[title={What you've built}]
One more word for giving reasons and opinions.
\end{grammarlens}

\subsection*{Your turn}
\begin{itemize}
\raggedright
  \item \emph{Say it:} \emph{rūpisu}
  \item \emph{Say it:} \emph{rūpisu}, once more
  \item \emph{Say it:} \emph{śikṣakaru hosa pāṭhagaḷannu rūpisidaru.} — The teachers devised new lessons.
  \item \emph{Recall:} say the Kannada for awareness, then the Kannada for diversity, then say \emph{rūpisu} again
\end{itemize}

\begin{tcolorbox}[breakable,colback=teal!4,colframe=teal!35!black,title={Before you move on}]
What does \kn{ರೂಪಿಸು} mean? (\textquotedblleft{}To devise; to shape\textquotedblright{}.) Say it once more.
\end{tcolorbox}
